\section{Términos clave: índice de palabras clave de este episodio}

\begin{longtable}{p{0.24\textwidth}p{0.34\textwidth}p{0.33\textwidth}}
\toprule
Término & Explicación breve & Papel en este episodio \\
\midrule
Robotics & Robótica: estudia la percepción, el control, la planificación y la ejecución física & Tema de todo el episodio. \\
Embodied AI & Inteligencia corporeizada: un agente percibe su entorno y actúa en él & Es un término común en el contexto chino, pero Tan Jie (谭杰) prefiere hablar de robots. \\
Sim2Real & Transferencia del entrenamiento en simulación al mundo real & Núcleo del primer cambio de paradigma de Tan Jie. \\
Reinforcement Learning & Aprendizaje por refuerzo: optimiza una política por ensayo y error mediante reward & Método importante para resolver la locomotion y el control. \\
PPO & Proximal Policy Optimization, algoritmo de aprendizaje por refuerzo de uso común & Uno de los antecedentes técnicos de los primeros trabajos de Tan Jie. \\
MPC & Model Predictive Control, control predictivo basado en modelos & Representa la línea tradicional de control de robots. \\
VLM & Vision-Language Model, modelo de visión y lenguaje & Puede ver imágenes y entender lenguaje, pero por lo general no produce acciones directamente. \\
VLA & Vision-Language-Action, modelo de visión, lenguaje y acción & Forma importante de los modelos fundacionales para robots. \\
Motion Transfer & Transferencia de movimiento: lleva las capacidades de movimiento a cuerpos robóticos distintos & Uno de los mecanismos clave de Gemini Robotics 1.5. \\
World Model & Modelo que predice cómo cambia el estado del mundo & Ayuda al robot a planificar las consecuencias de sus acciones. \\
V-L-V & Vision-Language-Vision & Una de las formas en que Tan Jie describe los modelos del mundo. \\
Synthetic Data & Datos sintéticos, generalmente obtenidos de simulación o generación por programa & Compensan la escasez de datos de robots reales. \\
Tactile Sensing & Percepción táctil: mide la fuerza de contacto, el deslizamiento y el material & Entrada importante para las manos diestras y la manipulación fina. \\
Dexterous Hand & Mano diestra: cuerpo robótico de varios dedos para operaciones complejas & Zona donde se concentran las dificultades del tacto y la manipulation. \\
Generalist & Robot de propósito general, capaz de generalizar entre tareas & Objetivo de largo plazo; terminará por desplazar al specialist. \\
Specialist & Robot especializado que resuelve tareas verticales específicas & Es más fácil de llevar a la práctica en el corto plazo. \\
\bottomrule
\end{longtable}

\subsection{Resumen de la sección}

Los términos de este episodio muestran que la robótica no es el nombre de un solo modelo, sino un conjunto de problemas de sistemas que abarcan la simulación, la acción, los modelos del mundo, el tacto, el hardware y la seguridad. Entender las relaciones entre los términos importa más que perseguir una demo aislada.
